% 摘要 fragment 不写 \section、\begin{abstract} 或 \end{abstract}。
% 用户内部规则：匿名题目、完整摘要、关键词必须同处一个逻辑页；正文从下一页开始。
% 题目优先单行，必要时最多两行且最终宽度近似均衡；关键词 3--6 个，每项整体不拆行。
% 若超页，删重复背景、次要过程和无决策价值的中间量；禁止缩字号、压行距、缩放或负间距。
% 通常按“总述—逐问—收束”组织，但自然段可按问题数量和逻辑合并/拆分，不固定为 2+N 段。
% 每个问题段压缩覆盖“任务—方法—结果—验证/风险—价值/去向”五个要素；这是信息要素，不是五段式硬排版。
% 所有数字、标签、比较结论均须可回溯到求解结果；没有目标真值时不得把预测写成准确率。
% 使用 \textbf{...} 选择性加粗主要模型、决定性数字（数值与单位一起）和核心结论/创新；普通年份、编号、中间量不加粗，禁止整句或整段加粗。

【总述，第 1 段】针对【对象】在【场景】中存在的【核心困难】（如数据缺失、分布偏移、约束耦合或不确定性）问题，本文构建【总体方法/证据链】。研究目标是【任务输出】，同时满足【关键边界或可信性要求】。

【对于问题一，第 2 段】对于问题一，依据【输入、机制或约束特征】将任务表述为【数学任务】，采用\textbf{【主要模型/算法】}求得\textbf{【带单位、精度或范围的关键结果】}；通过【理论互验、残差、基线对比、约束校验或扰动】确认\textbf{【核心结论或风险】}，并输出【供本问回答或后问使用的量】。

【对于问题二，第 3 段】对于问题二，在【数据隔离、时序、分组或资源条件】下采用\textbf{【主要模型/算法】}，并以【评价指标或硬约束】选择【方案】。相较【同口径基线】，核心指标由【基线值】变为\textbf{【本方案关键值】}；\textbf{【验证结论】}表明该结果在【适用范围】内成立，代价或限制为【运行成本/不确定性】。

【对于问题三至问题 N】每个问题各写一个自然段，均以“对于问题三，……”或“对于问题 N，……”开始。依次说明【输入和困难】、【方法】、【关键结果】和【验证/风险】；对无标签、不可验证或高风险部分，只输出【不确定性、置信度、风险等级或复核建议】，不夸大结论。

【总结，第 N+2 段】最后，通过【独立重算、交叉/滚动验证、消融、情景比较或约束校验】核验【核心结果】。本文的贡献在于【可被事实支撑的创新点】；研究结果可为【实际应用】提供【决策/分析】依据，但仅适用于【边界条件】，仍需由【后续数据或实验】进一步验证。
